% =========================================================
% DESPLIEGUE Y TRANSICIÓN (tarea O3-06) — fase de transición de RAD
% QUÉ VA: lo que pasó al llevar el sistema al centro. Es la condición previa de
%         la evaluación: sin uso estable no hay postest válido.
% CONTENIDO:
%   1. Despliegue: VPS (proveedor, vCPU, RAM, costo mensual), contenedores en
%      ejecución, dominio y certificado, respaldo diario hacia Google Drive de la
%      clínica. Contrastar con lo previsto en despliegue.tex (consumo estimado
%      0,6-1 GB): poner el consumo medido (docker stats).
%      Si se terminó en la instalación local, decirlo y por qué (alternativa
%      ya prevista en tab:ob2_local).
%   2. Carga inicial de datos: profesionales, horarios, catálogos (CIE-10,
%      tratamientos, fármacos) y pacientes. Decidir: ¿se migran los pacientes
%      del Excel o se registran al llegar? (afecta el tiempo de recepción en
%      las primeras semanas). Indicar volumen cargado.
%   3. Capacitación: fecha, duración y a quién (secretaria, odontólogos).
%   4. Periodo de estabilización: fechas de inicio de uso real y semanas de uso
%      antes del postest (recomendado: al menos 2 semanas). Incidencias y
%      ajustes realizados en ese periodo.
% EXTENSIÓN: 1 tabla de hitos (fecha | hito) + 1-2 párrafos.
% =========================================================

\subsubsection{Despliegue y transición}

\textit{[2--3 líneas: el sistema se desplegó según el esquema del diseño (Figura~\ref{fig:ob2_despliegue}); recursos del servidor y consumo medido frente a lo estimado; respaldo diario verificado con una restauración.]}

La Tabla~\ref{tab:ob3_hitos} resume la transición desde el despliegue hasta la medición postest.

\begin{table}[H]
\centering
\caption{Hitos de la transición}
\label{tab:ob3_hitos}
\fontsize{9}{11}\selectfont
\renewcommand{\arraystretch}{1.2}
\setlength{\tabcolsep}{4pt}
\begin{tabularx}{\linewidth}{|>{\raggedright\arraybackslash}p{2.6cm}|>{\raggedright\arraybackslash}X|}
\hline
\textbf{Fecha} & \textbf{Hito} \\ \hline
\textit{[dd/mm/2026]} & Despliegue en producción \\ \hline
\textit{[dd/mm/2026]} & Carga inicial: profesionales, horarios, catálogos y pacientes \\ \hline
\textit{[dd/mm/2026]} & Capacitación del personal \\ \hline
\textit{[dd/mm/2026]} & Inicio del uso real y del agente con el número de recepción \\ \hline
\textit{[dd/mm/2026]} & Fin del periodo de estabilización \\ \hline
\textit{[dd/mm/2026]} & Medición postest (tres jornadas) \\ \hline
\end{tabularx}
\end{table}

\textit{[2--3 líneas: incidencias del periodo de estabilización y ajustes realizados antes del postest.]}
